%% kapa: exam=1 round=36 year=2025 period=1 form=A subject=real_estate subject_ko=부동산학원론 seq=34 number=114
%% source: Q-Net 감정평가사 기출문제 (cst00309) artlSeq=5231525 file=raw/1008/5231525/2025년 제36회 감정평가사 1차시험 1교시 문제지 원본.pdf
%% conversion: kordoc 4.24.51 -> tools/gen1.py (auto); review=proofread-sonnet; flags=none
\begin{question}{114}
부동산권리분석활동을 위한 자료의 조사ㆍ확인 및 분석에 관한 설명으로 옳지 않은 것은?
\begin{choices}
  \item ｢공간정보의 구축 및 관리 등에 관한 법령｣상 지적도에 기재된 분석대상 부동산의 지목이 ‘공’ 으로 표기되어 있어, 지목을 공원으로 판단하였다.
  \item 구거는 용수(用水) 또는 배수(排水)를 위하여 일정한 형태를 갖춘 인공적인 수로ㆍ둑 및 그 부속시설물의 부지와 자연의 유수(流水)가 있거나 있을 것으로 예상되는 소규모 수로부지로 지적도에는 ‘구’ 로 표기한다.
  \item 유지(溜池)는 자연의 유수(流水)가 있거나 있을 것으로 예상되는 토지로 지적도에는 ‘유’ 로 표기한다.
  \item 건물의 소재지, 구조, 용도 등의 사실관계를 건축물대장을 통하여 확인ㆍ판단하였다.
  \item 권리분석보고서에는 대상부동산 및 의뢰인, 권리분석의 목적, 판단결과의 표시 및 이유, 권리분석의 방법 및 성격, 수집한 자료의 목록 및 면책사항 등이 포함된다.
\end{choices}
\end{question}
\answer{3}
